\section{Supplementary reading and proof boundaries}\label{sec:reading-audit}

The following evaluations explain why some promising references
were useful as reading companions but not chosen as the proof
of a stack axiom.

\begin{description}[style=nextline]
\item[Vakil and the Pixton course lead.]
Vakil \cite[Sections 2.4--2.7, pp.~9--11]{Vakil} gives geometric
motivation for stable curves and their compactification.
Pixton's \href{https://public.websites.umich.edu/~pixton/classes/697-F22/index.html}{Math 697 course reference page}
also directs readers to a Vakil article for an informal introduction.
This is useful orientation, not a proof of object descent.
Vakil's later \href{https://math.stanford.edu/~vakil/22-245C/22-245Cnotes2022-05-17.pdf}{2022 course notes},
Section 45.6, p.~88, give another stable-map algebraicity route
using a mapping stack, but leave part of its representability
discussion as an exercise. The direct family and Hilbert methods
fit the present scope better.

\item[Zvonkine.]
The stable-curve discussion in
\cite[Section 1.4, pp.~14--18]{Zvonkine} is particularly good
for examples of compactification and special-point stability.
Its account is intentionally informal about foundations, so
it does not replace the scheme descent proof.

\item[Student expositions.]
Belfiori \cite[Section 1.1, pp.~2--5]{Belfiori} explicitly
describes cartesian family arrows and the role of a canonical
polarization in descent, citing Alper and Vistoli. Its theorem
about the stable compactification's geometric properties is
stated rather than proved. Monavari
\cite[Sections 1.2--1.3, pp.~11--16]{Monavari} gives additional
practice with fibered categories and Deligne--Mumford stacks;
it does not prove both of our complete constructions.
These were read as supplementary expositions, not as original
sources for the classical theorems.

Niven Achenjang's \href{https://people.math.harvard.edu/~achenjang/assets/pdf/UW__Math_582C_Notes.pdf}{\emph{Introduction to Stacks and Moduli Notes}}
for Alper's Winter 2021 course, pp.~19--21 and Sections 5.2
and 5.6, pp.~24 and 28--29, gives a conversational companion
to descent and Hilbert presentations. We used the checked
Alper text for the formal results. Cavalieri's
\href{https://www.math.colostate.edu/~renzo/teaching/Moduli16/Fields.pdf}{\emph{Moduli Spaces of Pointed Rational Curves}},
Fields graduate summer school, July 2016, Section 1.2,
pp.~10--11, discusses families and their equivalence; its
genus-zero focus does not supply the two general stack proofs.

\item[Other leads and inaccessible material.]
The search also followed the suggested Cavalieri, Caporaso,
Farkas, Graber, Pandharipande, Pixton, and Ross leads.
Several results concerned intersection theory, particular
components, or later applications of moduli stacks, which
does not establish the foundational gluing claim. For example,
Farkas's \emph{Regular components of moduli spaces of stable
maps} directs readers to Fulton--Pandharipande for construction.
A search-indexed Imperial lecture-note file advertised a
stable-map algebraicity section but the actual file could not
be retrieved; no claim here relies on that unseen section.
The UCLA student stack seminar supplied another reading list,
but the handwritten PDFs could not be inspected through the
available web reader and were not used as evidence.
\end{description}

\subsection{What has and has not been reconstructed}

The two stack proofs are complete relative to their stated
scheme-theoretic inputs: morphism descent, duality and base
change, ampleness criteria, and polarized effective descent.
The hypotheses of those inputs are checked at their uses.
For the later properties, Hilbert representability and bounded
embeddings, relevant infinitesimal lifting results, criteria
for Deligne--Mumford stacks and properness, and the stable
reduction/extension theorems are substantial black boxes.
The guided proofs explain how they imply the propositions;
they do not replace the full proofs of those external theorems.

We did not find one accessible pedagogical source spelling out
all the requested pointed-curve and stable-map descent checks
in Fantechi's notation. The present chapters supply that
reconstruction from the verified general descent theorems.
Behrend--Manin's reference to an earlier preliminary properness
proof was not independently traced to an accessible text;
Fulton--Pandharipande supplies a verifiable later family-level
argument. No historical-priority assertion is made for that
unverified preliminary source.

There is no unresolved classical stack claim in the stated
setup once the declared external results are accepted. There
is also no claim to have constructed the intended cuspidal
moduli problem: its defining stability condition, descent
polarization, boundedness, diagonal, and limit theorems are
the separate research tasks in Chapter~\ref{ch:cusps}.
